\scanpage{183}\mspage{39}
% archive sheet number: 183

\margin{Ex : les Cat-fibrations}

$(3.8.)\quad$\raisebox{-5.5\baselineskip}[0pt][0pt]{$\left\{\vphantom{\rule[-5.5\baselineskip]{0pt}{11\baselineskip}}\right.$}\,$W_0$-lisse ($\Longleftrightarrow$ les inclusions $I'_i\unsure{\longrightarrow|\mathbin{-}}{}_{i\backslash}I'$ sont $W_0$-asphériques, i.e. les fibres des flèches localisées \deleted{$\unsure{A}$} $I'_{/i'}\unsure{\longrightarrow|\mathbin{-}}I_{/i}$ $(i=ui')$ sont $0$-connexes

$W_0$-parfait $=W_0$-propre et $W_0$-lisse

$W_0$-équivalences universelles ($\pi_0$ bijectif et le reste pour tout changement de base.)

\margin{Ex : les \deleted{\unsure{flèches}} $W_0$-lisses ou $W_0$-propres, à fibres $0$-connexes}

$\Downarrow$ impl. \uuline{stricte} (et comment)

\inserted{Ce qui impli\unsure{que} les \deleted{$W$}-fibrations (au \unsure{sens} faible) à fibres asphériques i.e. $0$-connexes}

$W_0$-fibration : $u^*:\mathrm{Cat}_{/I}\longrightarrow\mathrm{Cat}_{/I'}$ transforme $W_0$-équivalences en $W_0$-équivalences

\margin{\deleted{Ex : les \unsure{fibrations} \unsure{propres}}}

$\Downarrow$ impl. \uuline{stricte} (et comment !)

\inserted{$\longleftrightarrow$ mais si les fibres sont $0$-connexes, \unsure{équivalentes} et signifie \unsure{les changements|les équivalences} de \unsure{base} des $W_0$-\unsure{fibrations} sont des $W_0$-\unsure{équiv.} universelles\ldots}

\uline{$W_0$-fibrat\deleted{\unsure{ion}} faible} ($\Longleftrightarrow$ \deleted{pour tout} $u^*:\mathrm{Cat}_{/I}\longrightarrow\mathrm{Cat}_{/I'}$ transforme homotopismes $J_1\overset{\deleted{\unsure{\sim}}}{\longrightarrow}J_2$ en $W_0$-équivalences -- et il suffit de l'exiger pour l'inclusion d'un objet final \inserted{(ou, au choix, initial)} \deleted{\unsure{ou pour l'in}} et pour l'inclusion d'un point dans $\Delta^1$)

\margin{\uline{Ex} les fl. $W_0$-\deleted{\unsure{lisses}} telles que $\pi_0(I'/I)=(i\mapsto\pi_0({}_{i\backslash}I'))=u_!(e_{I'{}^\wedge})$ soit loc. constant sur $I$ i.e. $\underline{\pi}_0(I'/I)$ soit un \uuline{rev.} étale de $I$, ou si $u^{\mathrm{opp}}$ satisfait à cette condition}

Signalons, comme \inserted{\unsure{candide|corollaire}} utilisation de ces notions, que pour un diagramme cartésien dans $\mathrm{Cat}$

\diagram{$(3.9.)\qquad$\begin{tikzcd}[ampersand replacement=\&]
I' \arrow[d,"u"'] \& J' \arrow[l,"{p'}"'] \arrow[d,"v"]\\
I \& J \arrow[l,"p"]
\end{tikzcd}$,$}

la flèche de chgt de base pour $u_*,v_*$ \inserted{pour les préfaisceaux} est \unsure{is} si $u$ est \inserted{$W_0$-}propre, ou $p$ $W_0$-lisse, et

la flèche de \unsure{chgt.} de base pour $u_!,v_!$ est \unsure{id|iso} si $u$ est $W_0$-\deleted{\unsure{faible}}\unsure{lisse} ou $p$ $W_0$-propre.

% slanted
\margin{Sauf erreur, \unsure{cela se démontre|cela découle} \uuline{\unsure{sans|d'une}} \uline{\unsure{caractérisation|considération}} \unsure{très} \unsure{simple} des morphismes $W_0$-propres, $W_0$-lisses \unsure{uniquement|impliquant}}

Je signale à ce propos que ces énoncés, qui à priori se \deleted{\unsure{démontraient}} formulent pour les préfaisceaux à valeurs dans $\mathrm{Ens}$, \unsure{restent|restent-ils} valables
